\section{Introduction}\label{sec:introduction}
For Schwartz functions $f_1,f_2,f_3:\R\to\C$, set
\[
T(f_1,f_2,f_3)(x)
 =\operatorname{p.v.}\int_{\R}
 f_1(x-t)f_2(x-2t)f_3(x-3t)\,\frac{dt}{t}.
\]
The principal value removes a symmetric interval about the origin.
Our main result is the following estimate.

\begin{theorem}\label{thm:main}
There is a finite constant $C$ such that
\[
 \norm{T(f_1,f_2,f_3)}_{L^1(\R)}
 \le C\prod_{j=1}^3\norm{f_j}_{L^3(\R)}
 \qquad(f_1,f_2,f_3\in\mathcal S(\R)).
\]
Consequently $T$ has a unique bounded trilinear extension from
$L^3(\R)\times L^3(\R)\times L^3(\R)$ to $L^1(\R)$.
\end{theorem}

Theorem~\ref{thm:main} resolves affirmatively the
$L^3\times L^3\times L^3\to L^1$ case of the standard trilinear
Hilbert transform conjecture for the fixed slopes $1,2,3$.
The full conjecture asks for a broader range of exponents; see
\cite[Conjecture~1.7]{HuLie2023}.
The result here concerns the displayed fixed operator.

The bilinear predecessor was treated by Lacey and Thiele~\cite{LT1997,LT1999}.
The difficulty in the trilinear setting is the interaction of cancellation across scales with
quadratic modulation.  For example, if
$c=(-1,3,-3,1)$, then
\[
 \sum_{j=0}^3c_j(x+jt)^a=0\qquad(a=0,1,2).
\]
Thus matched quadratic phases can persist in the associated four-linear
form.  Quantitative information at one scale must be organized so that
this persistence does not produce a loss proportional to the number
of scales.
Tao~\cite{Tao2015} obtained cancellation for truncated multilinear
Hilbert transforms, improving the elementary logarithmic truncation
bound to a sublogarithmic one.  That argument established a connection
with higher-order uniformity.  The straight-line trilinear boundedness
problem remained open in the account of Hu and Lie~\cite[Section~1.2]{HuLie2025}.
Their results for curved transforms~\cite{HuLie2023,HuLie2025} address
a different source of oscillation.

The inverse theorem of Leng, Sah and Sawhney~\cite[Theorem~1.2]{LSS2024}
provides the quantitative structural input used below.  We use its
degree-two case and prove the continuous chart and localization
statements needed for the singular integral.  The inverse theorem alone
does not sum the scales.  The proof of that summation has three main
features.

First, private orthogonal coordinates make the coefficients of each
selected column recoverable.  Lemma~\ref{lem:compression} then controls
products of inherited projections uniformly in spatial depth.  Its loss
is a power of the logarithm of a common budget for column sizes, inverse
private-coordinate sizes, and the number of available labels along a
spatial path; the total number of labels on disjoint branches is unrestricted.

Second, we localize curvature jointly in all four arguments and retain
the separation needed when projection adjoints are evaluated on earlier
cells.  Transferring complementary filters leaves coupled narrow
bands.  The comparison in Lemma~\ref{lem:short-block} gives a
power of block length whose exponent tends to zero as the auxiliary
exponent decreases to two.  This control is essential to the final
absorption.

Third, each matched column receives one fixed mode expansion.
Frequency graphs separate errors with square-summable energy from
coherent components with hereditary representatives.
Lemma~\ref{lem:finite-sparse-count} bounds the total mass of progression
tests on scales of small density, with a positive-power saving in the
density.  Its application uses all four anchor cancellations and
polynomial dimension costs.  The compression and counting lemmas isolate
tools for other progression forms with the same projection and
cancellation structures.

\subsection*{Organization and notation}
Section~\ref{sec:outer} reduces the operator estimate to a uniform
maximal-function bound for local forms.  Sections~\ref{sec:compression}
and~\ref{sec:charts} develop the projection and quadratic tools.
Section~\ref{sec:base} produces a pair-small base form with a bounded
number of active scales on each path.  Section~\ref{sec:localization}
localizes this form by curvature.  Section~\ref{sec:sides} gives the side
estimates under a square-energy hypothesis, which is proved in
Section~\ref{sec:linearization}.  Section~\ref{sec:counting} estimates the
remaining core terms.  Section~\ref{sec:completion} chooses the parameters
and completes the absorption.

We write $e(t)=\exp(2\pi i t)$ and use progression coordinates
$y_j=x+jt$, $0\le j\le3$.
Inner products are linear in the first argument.  Replacing the shifts
$x-jt$ by $x+jt$ in the principal-value integral changes the operator's
sign, by the substitution $t\mapsto-t$.
Norms and inner products carrying an interval
subscript use normalized Lebesgue measure on that interval.  Other
function norms use ordinary Lebesgue measure.  In vector-valued
arguments the scalar component entering the form is called the physical
component.  All depth trees used in the estimates are finite; constants
will be uniform in their depth.  The polynomial and quasipolynomial
budget conventions are specified in Section~\ref{sec:compression}.
